% !TEX root = ../thesis.tex
\chapter{Motivácia}\label{ch:motivation}

Rastúca zložitosť moderných webových systémov spolu s požiadavkami na krátke odozvy, vysokú dostupnosť a flexibilné rozširovanie funkcií vytvára tlak na architektonické rozhodnutia v každom softvérovom projekte. Vývojové tímy často stoja pred dilemou \cite{Newman2021,Richardson2018}, či uprednostniť jednoduchý \glslink{monolith}{monolitický prístup}, ktorý zrýchľuje prototypovanie a znižuje prevádzkové nároky, alebo implementovať \glslink{microservice}{mikroservisnú architektúru}, ktorá prináša nezávislé škálovanie, distribuované spracovanie a možnosť samostatného nasadzovania komponentov \cite{FowlerMicroservices}. Táto práca vychádza z potreby poskytnúť empiricky podložené odporúčania založené na reálnej implementácii a porovnaní dvoch architektonicky odlišných riešení \cite{FowlerMonolithFirst}.

Vo všeobecnosti platí, že výber architektúry má zásadný vplyv na celý životný cyklus softvéru – od návrhu cez implementáciu až po prevádzku a dlhodobú údržbu. V malých a stredne veľkých projektoch býva \glslink{monolith}{monolitický prístup} často rýchlejší na vývoj a umožňuje koncentrovať všetku biznis logiku na jednom mieste. S rastúcimi požiadavkami na modularitu a paralelný vývoj však môže byť monolit bariérou. \glslink{microservice}{Mikroservisy} naopak umožňujú škálovanie jednotlivých častí systému, jednoduchšie testovanie izolovaných modulov a flexibilné nasadzovanie. Ich výhodou je aj odolnosť voči zlyhaniam – porucha jednej služby neohrozí celok. Na druhej strane prinášajú zvýšenú prevádzkovú komplexitu, nároky na monitoring, logovanie a správu distribuovaných dát.

Hru Dáma som zvolil za modelovú doménu z niekoľkých dôvodov: pravidlá sú jednoznačné, rozsah stavu je obmedzený a možné scenáre použitia dobre reprezentujú bežné požiadavky webových aplikácií (interaktívne aktualizácie stavu, perzistencia výsledkov, správa používateľov, zobrazenie štatistík). Navyše ide o typickú aplikáciu so zreteľne oddeliteľnými doménami – hráči, hra, skóre, komentáre, hodnotenia – čo vytvára ideálne podmienky pre demonštráciu rozdielov medzi \glslink{monolith}{monolitom} a \glslink{microservice}{mikroservismi}. Motivácia je teda kombináciou pedagogického zámeru (ukázať rozdiely v praxi), akademického prínosu (zmapovať vplyv architektúry na správanie systému) a praktického zamerania (poskytnúť odporúčania pre reálne projekty).

Zároveň je rastúci dopyt po distribuovaných systémoch a cloud-native riešeniach trendom posledného desaťročia. Podnikové systémy, e-commerce platformy či streamingové služby sa presúvajú od \glslink{monolith}{monolitov} k \glslink{microservice}{mikroservisom} práve kvôli potrebe škálovania, resiliencie a rýchlej inovácie. Táto práca umožňuje pochopiť, do akej miery sú tieto výhody relevantné aj pri menšej aplikácii s jasne definovanými funkčnými hranicami.

Okrem technických a architektonických hľadísk zohráva dôležitú úlohu aj praktický rozmer práce v kontexte výučby a samostatného štúdia softvérového inžinierstva. Študenti sa často stretávajú s pojmami ako „\glslink{monolith}{monolit}“, „\glslink{microservice}{mikroservis}“, „škálovateľnosť“ či „observabilita“ len na teoretickej úrovni, bez priameho prepojenia na konkrétny projekt. Vytvorenie dvoch reálnych implementácií tej istej aplikácie – jednej \glslink{monolith}{monolitickej} a jednej založenej na \glslink{microservice}{mikroservisoch} – umožňuje tieto abstraktné pojmy konkretizovať, porovnať ich v praxi a lepšie pochopiť, aké kompromisy jednotlivé prístupy prinášajú. Táto práca tak môže slúžiť aj ako študijný materiál pre ďalších študentov, ktorí sa chcú oboznámiť s modernými architektonickými vzormi na zrozumiteľnom a dobre ohraničenom príklade.